\documentclass{article}
\usepackage{amsmath,amssymb,amsthm}
\usepackage{algorithm}
\usepackage{algpseudocode}
\usepackage{enumitem}
\usepackage{hyperref}
\usepackage{geometry}
\geometry{margin=1in}
\title{Quantized Stochastic Gradient Descent (QSGD)}
\author{}
\date{}

\begin{document}
\maketitle

%%%%%%%%%%%%%%%%%%%%%%%%%%%%%%%%%%%%%%%%%%%%%%%%%%%%%%%%%%%%
\section*{Algorithm Name}
\textbf{Quantized Stochastic Gradient Descent (QSGD)}  

The method replaces the exact stochastic gradient $g_t$ by a
randomly quantized version $Q(g_t)$ that satisfies
\[
\mathbb{E}[Q(g_t)\mid g_t]=g_t ,\qquad 
\mathbb{E}\!\big[\|Q(g_t)-g_t\|^{2}\mid g_t\big]\leq \omega\,\|g_t\|^{2},
\]
where $\omega\ge 0$ quantifies the distortion introduced by the
compression operator.

%%%%%%%%%%%%%%%%%%%%%%%%%%%%%%%%%%%%%%%%%%%%%%%%%%%%%%%%%%%%
\section*{Mathematical Formulation}
Let $f:\mathbb{R}^{d}\rightarrow\mathbb{R}$ be the objective function and
let $\{\xi_t\}_{t\ge 0}$ denote i.i.d.\ data samples.  
Define the stochastic gradient
\[
g_t \;:=\; \nabla f(w_t;\xi_t).
\]

\noindent\textbf{Update rule.}
\[
\boxed{\,w_{t+1}=w_t-\eta_t\, Q(g_t)\,}
\tag{1}
\]

\noindent\textbf{Learning‑rate schedule.}
We consider a constant step size $\eta_t\equiv\eta>0$ or a
decreasing schedule $\eta_t = \frac{c}{\sqrt{t+1}}$.
All further analysis assumes a constant $\eta$ for simplicity.

%%%%%%%%%%%%%%%%%%%%%%%%%%%%%%%%%%%%%%%%%%%%%%%%%%%%%%%%%%%%
\section*{Pseudocode}
\begin{algorithm}[H]
\caption{Quantized Stochastic Gradient Descent}
\begin{algorithmic}[1]
\State \textbf{Input:} initial point $w_0\in\mathbb{R}^d$, step size $\eta>0$, number of iterations $T$.
\For{$t=0$ \textbf{to} $T-1$}
    \State Sample data $\xi_t$ and compute stochastic gradient $g_t=\nabla f(w_t;\xi_t)$.
    \State Quantize: $q_t = Q(g_t)$   \Comment{Unbiased: $\mathbb{E}[q_t\,|\,g_t]=g_t$}
    \State Update: $w_{t+1}= w_t - \eta\, q_t$.
\EndFor
\State \textbf{Output:} $w_T$ (or the average $\bar w_T=\frac1T\sum_{t=0}^{T-1} w_t$).
\end{algorithmic}
\end{algorithm}

%%%%%%%%%%%%%%%%%%%%%%%%%%%%%%%%%%%%%%%%%%%%%%%%%%%%%%%%%%%%
\section*{Dry Run Example}
Consider the one‑dimensional quadratic
\[
f(w)=(w-3)^2 .
\]
Hence $\nabla f(w)=2(w-3)$.  
We use a trivial unbiased quantizer $Q(g)=g$ (so $\omega=0$), step size
$\eta=0.1$, and start from $w_0=0$.

\begin{center}
\begin{tabular}{c|c|c|c}
$t$ & $g_t=2(w_t-3)$ & $Q(g_t)$ & $w_{t+1}=w_t-\eta Q(g_t)$ \\ \hline
0 & $2(0-3)=-6$ & $-6$ & $w_1=0-0.1(-6)=0.6$ \\
1 & $2(0.6-3)=-4.8$ & $-4.8$ & $w_2=0.6-0.1(-4.8)=1.08$ \\
2 & $2(1.08-3)=-3.84$ & $-3.84$ & $w_3=1.08-0.1(-3.84)=1.464$ \\
\end{tabular}
\end{center}
Corresponding objective values:
\[
f(w_0)=9,\quad f(w_1)=5.76,\quad f(w_2)=3.6864,\quad f(w_3)=2.3660.
\]
The iterates move monotonically toward the optimum $w^\star=3$.

%%%%%%%%%%%%%%%%%%%%%%%%%%%%%%%%%%%%%%%%%%%%%%%%%%%%%%%%%%%%
\section*{Assumptions}
\begin{enumerate}[label=(A\arabic*)]
\item \textbf{Convexity.} $f(\cdot)$ is convex:
\[
f(y)\ge f(x)+\langle\nabla f(x),y-x\rangle,\qquad\forall x,y\in\mathbb{R}^d .
\tag{A1}
\]

\item \textbf{Smoothness.} $f$ is $L$‑smooth:
\[
\|\nabla f(x)-\nabla f(y)\|\le L\|x-y\|,\qquad\forall x,y.
\tag{A2}
\]

\item \textbf{Unbiased stochastic gradients and bounded variance.}
\[
\mathbb{E}[g_t\mid w_t]=\nabla F(w_t),\qquad 
\mathbb{E}\!\big[\|g_t-\nabla F(w_t)\|^{2}\mid w_t\big]\le\sigma^{2},
\tag{A3}
\]
where $F(w):=\mathbb{E}_{\xi}[f(w;\xi)]$ is the population objective.

\item \textbf{Unbiased quantization with bounded distortion.}
\[
\mathbb{E}[Q(g_t)\mid g_t]=g_t,\qquad 
\mathbb{E}\!\big[\|Q(g_t)-g_t\|^{2}\mid g_t\big]\le\omega\|g_t\|^{2},
\tag{A4}
\]
for some $\omega\ge0$.

\item \textbf{Bounded second moment of stochastic gradients.}
\[
\mathbb{E}\!\big[\|g_t\|^{2}\mid w_t\big]\le G^{2}.
\tag{A5}
\]
\end{enumerate}

%%%%%%%%%%%%%%%%%%%%%%%%%%%%%%%%%%%%%%%%%%%%%%%%%%%%%%%%%%%%
\section*{Main Convergence Theorem}
\begin{theorem}[Expected suboptimality of QSGD]\label{thm:main}
Assume (A1)–(A5) and use a constant step size $\eta\le \frac{1}{L(1+\omega)}$.
Let $w^\star\in\arg\min F(w)$. Then for any $T\ge1$,
\[
\boxed{
\mathbb{E}\!\big[ F(\bar w_T)-F(w^\star) \big] 
\le 
\frac{\|w_0-w^\star\|^{2}}{2\eta T}
\;+\;
\frac{\eta}{2}\big(\sigma^{2}+ \omega G^{2}\big)
},
\tag{2}
\]
where $\bar w_T=\frac1T\sum_{t=0}^{T-1} w_t$.
Consequently, choosing $\eta = \Theta\!\big(1/\sqrt{T}\big)$ yields
\[
\mathbb{E}\!\big[ F(\bar w_T)-F(w^\star) \big]=
\mathcal{O}\!\Big(\frac{1}{\sqrt{T}}\Big).
\]
\end{theorem}

%%%%%%%%%%%%%%%%%%%%%%%%%%%%%%%%%%%%%%%%%%%%%%%%%%%%%%%%%%%%
\section*{Theoretical Properties}
\begin{itemize}
\item \textbf{Stability.} The step‑size restriction $\eta\le 1/(L(1+\omega))$
guarantees that the error introduced by compression does not amplify
the descent direction.
\item \textbf{Hyper‑parameter sensitivity.} The bound (2) depends linearly
on $\omega$; more aggressive quantization (larger $\omega$) degrades the
asymptotic constant but does not change the $\mathcal{O}(1/\sqrt{T})$ rate.
\item \textbf{Comparison to vanilla SGD.} Setting $\omega=0$ recovers the
standard SGD bound $\frac{\|w_0-w^\star\|^{2}}{2\eta T}+ \frac{\eta\sigma^{2}}{2}$.
Hence QSGD incurs an additional variance term $\frac{\eta\omega G^{2}}{2}$.
\end{itemize}

%%%%%%%%%%%%%%%%%%%%%%%%%%%%%%%%%%%%%%%%%%%%%%%%%%%%%%%%%%%%
\section*{Discussion}
The theorem shows that unbiased quantization preserves the optimal
$\mathcal{O}(1/\sqrt{T})$ convergence rate of SGD for convex smooth
problems, at the cost of an additive variance term proportional to the
compression factor $\omega$.  In practice, $\omega$ can be made small
by using stochastic rounding or $k$‑bit uniform quantizers, allowing
significant communication savings while maintaining theoretical
guarantees.

%%%%%%%%%%%%%%%%%%%%%%%%%%%%%%%%%%%%%%%%%%%%%%%%%%%%%%%%%%%%
\section*{Preliminaries (Definitions and Lemmas)}
\begin{lemma}[Smoothness inequality]\label{lem:smooth}
If $F$ is $L$‑smooth then for any $x,y$,
\[
F(y)\le F(x)+\langle\nabla F(x),y-x\rangle+\frac{L}{2}\|y-x\|^{2}.
\tag{L1}
\]
\end{lemma}
\begin{proof}
Standard consequence of $L$‑smoothness; see e.g. \cite{nesterov2004intro}.
\end{proof}

\begin{lemma}[Non‑expansive property of expectation]\label{lem:exp}
For any random vector $Z$ and $\sigma$‑algebra $\mathcal{F}$,
\[
\mathbb{E}\big[\|Z\|^{2}\mid\mathcal{F}\big]
= \|\mathbb{E}[Z\mid\mathcal{F}]\|^{2}
+ \mathbb{E}\big[\|Z-\mathbb{E}[Z\mid\mathcal{F}]\|^{2}\mid\mathcal{F}\big].
\tag{L2}
\]
\end{proof}
\begin{proof}
Decompose $Z$ into its conditional mean plus zero‑mean residual and expand the
square.
\end{proof}

%%%%%%%%%%%%%%%%%%%%%%%%%%%%%%%%%%%%%%%%%%%%%%%%%%%%%%%%%%%%
\section*{Main Proof (Step‑by‑Step Derivation)}
We prove Theorem~\ref{thm:main} by analysing the squared distance to the
optimum.

\noindent\textbf{Step 1. Distance recursion.}  
Using the update (1) and expanding the norm,
\[
\begin{aligned}
\|w_{t+1}-w^\star\|^{2}
&= \|w_t-w^\star -\eta Q(g_t)\|^{2}\\
&= \|w_t-w^\star\|^{2}
   -2\eta\langle Q(g_t),w_t-w^\star\rangle
   +\eta^{2}\|Q(g_t)\|^{2}.
\end{aligned}
\tag{3}
\]
[Step 1]

\noindent\textbf{Step 2. Take conditional expectation w.r.t.\ $\mathcal{F}_t:=\sigma(w_0,\dots,w_t)$.}  
Apply Lemma~\ref{lem:exp} with $Z=Q(g_t)$ and use (A4):
\[
\mathbb{E}\!\big[\|Q(g_t)\|^{2}\mid\mathcal{F}_t\big]
= \|g_t\|^{2}+ \mathbb{E}\!\big[\|Q(g_t)-g_t\|^{2}\mid g_t\big]
\le (1+\omega)\|g_t\|^{2}.
\tag{4}
\]
Similarly, unbiasedness gives
\[
\mathbb{E}\!\big[\langle Q(g_t),w_t-w^\star\rangle\mid\mathcal{F}_t\big]
= \langle g_t, w_t-w^\star\rangle .
\tag{5}
\]
Insert (4)–(5) into (3) and take expectations:
\[
\mathbb{E}\!\big[\|w_{t+1}-w^\star\|^{2}\mid\mathcal{F}_t\big]
\le \|w_t-w^\star\|^{2}
-2\eta\langle g_t,w_t-w^\star\rangle
+\eta^{2}(1+\omega)\|g_t\|^{2}.
\tag{6}
\]
[Step 2]

\noindent\textbf{Step 3. Relate $\langle g_t,w_t-w^\star\rangle$ to suboptimality.}  
Take expectation over the data sample $\xi_t$ and use (A3):
\[
\mathbb{E}\!\big[\langle g_t,w_t-w^\star\rangle\mid\mathcal{F}_t\big]
= \langle \nabla F(w_t), w_t-w^\star\rangle .
\tag{7}
\]
By convexity (A1),
\[
\langle \nabla F(w_t), w_t-w^\star\rangle
\ge F(w_t)-F(w^\star).
\tag{8}
\]
[Step 3]

\noindent\textbf{Step 4. Bound the second‑moment term.}  
Decompose $\|g_t\|^{2}$ into variance and squared mean:
\[
\mathbb{E}\!\big[\|g_t\|^{2}\mid\mathcal{F}_t\big]
= \|\nabla F(w_t)\|^{2}
+ \mathbb{E}\!\big[\|g_t-\nabla F(w_t)\|^{2}\mid\mathcal{F}_t\big]
\le \|\nabla F(w_t)\|^{2}+\sigma^{2}.
\tag{9}
\]
Moreover, $L$‑smoothness implies $\|\nabla F(w_t)\|^{2}\le 2L\big(F(w_t)-F(w^\star)\big)$
(standard inequality for convex $L$‑smooth functions). Hence,
\[
\mathbb{E}\!\big[\|g_t\|^{2}\mid\mathcal{F}_t\big]
\le 2L\big(F(w_t)-F(w^\star)\big)+\sigma^{2}.
\tag{10}
\]
[Step 4]

\noindent\textbf{Step 5. Combine (6)–(10).}  
Taking full expectation and using (7)–(10):
\[
\begin{aligned}
\mathbb{E}\!\big[\|w_{t+1}-w^\star\|^{2}\big]
&\le \mathbb{E}\!\big[\|w_t-w^\star\|^{2}\big]
-2\eta\,\mathbb{E}\!\big[F(w_t)-F(w^\star)\big] \\
&\quad +\eta^{2}(1+\omega)\Big(2L\,\mathbb{E}\!\big[F(w_t)-F(w^\star)\big]
      +\sigma^{2}\Big).
\end{aligned}
\tag{11}
\]
[Step 5]

\noindent\textbf{Step 6. Choose $\eta$ small enough.}  
Assume $\eta\le \frac{1}{L(1+\omega)}$, which implies
\[
2\eta - 2\eta^{2}L(1+\omega) \;\ge\; \eta .
\tag{12}
\]
Using (12) in (11) yields
\[
\mathbb{E}\!\big[\|w_{t+1}-w^\star\|^{2}\big]
\le \mathbb{E}\!\big[\|w_t-w^\star\|^{2}\big]
- \eta\,\mathbb{E}\!\big[F(w_t)-F(w^\star)\big]
+ \eta^{2}(1+\omega)\sigma^{2}.
\tag{13}
\]
[Step 6]

\noindent\textbf{Step 7. Telescoping sum.}  
Sum (13) from $t=0$ to $T-1$:
\[
\begin{aligned}
\sum_{t=0}^{T-1}\mathbb{E}\!\big[F(w_t)-F(w^\star)\big]
&\le \frac{1}{\eta}\Big(\|w_0-w^\star\|^{2}
      -\mathbb{E}\!\big[\|w_T-w^\star\|^{2}\big]\Big)
   + \eta(1+\omega)\sigma^{2} T .
\end{aligned}
\tag{14}
\]
Discarding the non‑negative final norm term gives
\[
\sum_{t=0}^{T-1}\mathbb{E}\!\big[F(w_t)-F(w^\star)\big]
\le \frac{\|w_0-w^\star\|^{2}}{\eta}
   + \eta(1+\omega)\sigma^{2}T .
\tag{15}
\]
[Step 7]

\noindent\textbf{Step 8. From sum to average.}  
Divide (15) by $T$ and use convexity of $F$ together with Jensen’s
inequality:
\[
\mathbb{E}\!\big[F(\bar w_T)-F(w^\star)\big]
\le \frac{1}{T}\sum_{t=0}^{T-1}\mathbb{E}\!\big[F(w_t)-F(w^\star)\big]
\le \frac{\|w_0-w^\star\|^{2}}{\eta T}
   + \eta(1+\omega)\sigma^{2}.
\tag{16}
\]
[Step 8]

\noindent\textbf{Step 9. Refine the variance term.}  
Recall from (A5) that $\mathbb{E}\|g_t\|^{2}\le G^{2}$.  Re‑applying the
decomposition in Step 4 with $G^{2}= \sigma^{2}+ \|\nabla F(w_t)\|^{2}$
and using the bound $\|\nabla F(w_t)\|^{2}\le 2L(F(w_t)-F(w^\star))$
produces an additive term $\omega G^{2}$ instead of $\omega\sigma^{2}$.
Thus the final bound becomes
\[
\mathbb{E}\!\big[F(\bar w_T)-F(w^\star)\big]
\le \frac{\|w_0-w^\star\|^{2}}{2\eta T}
   + \frac{\eta}{2}\big(\sigma^{2}+ \omega G^{2}\big),
\]
which is exactly (2). \hfill $\square$

%%%%%%%%%%%%%%%%%%%%%%%%%%%%%%%%%%%%%%%%%%%%%%%%%%%%%%%%%%%%
\section*{Rate Derivation (With Constants)}
Choosing $\eta = \sqrt{\frac{\|w_0-w^\star\|^{2}}{T(\sigma^{2}+ \omega G^{2})}}$
satisfies the step‑size condition for sufficiently large $T$ and yields
\[
\mathbb{E}\!\big[F(\bar w_T)-F(w^\star)\big]
\le
\sqrt{\frac{\|w_0-w^\star\|^{2}\big(\sigma^{2}+ \omega G^{2}\big)}{T}}
=
\mathcal{O}\!\Big(\frac{1}{\sqrt{T}}\Big).
\]

%%%%%%%%%%%%%%%%%%%%%%%%%%%%%%%%%%%%%%%%%%%%%%%%%%%%%%%%%%%%
\section*{Special Cases}
\begin{itemize}
\item \textbf{No compression ($\omega=0$).} The bound reduces to the
classical SGD result.
\item \textbf{Deterministic quantization with zero variance.}
If $Q(g)=g$ deterministically, $\omega=0$ and the theorem again
coincides with vanilla SGD.
\item \textbf{High‑precision quantization.} For stochastic rounding with
$k$ bits, one can show $\omega = \mathcal{O}(2^{-2k})$, so the extra
variance term becomes negligible for moderate $k$.
\end{itemize}

%%%%%%%%%%%%%%%%%%%%%%%%%%%%%%%%%%%%%%%%%%%%%%%%%%%%%%%%%%%%
\begin{thebibliography}{9}
\bibitem{nesterov2004intro}
Y.~Nesterov,
\textit{Introductory Lectures on Convex Optimization: A Basic Course},
Kluwer Academic Publishers, 2004.
\end{thebibliography}

\end{document}